\documentclass[11pt]{article}

\usepackage[margin=1in]{geometry}
\usepackage{amsmath,amssymb,amsthm}
\usepackage{bm}
\usepackage{xspace}
\usepackage{booktabs}
\usepackage{array}
\usepackage{enumitem}
\usepackage{listings}
\usepackage{xcolor}
\usepackage{hyperref}
\usepackage[capitalize]{cleveref}
\hypersetup{colorlinks=true, linkcolor=blue, urlcolor=blue, citecolor=blue}

% ---- Shorthands ----
\newcommand{\skewt}{skew-$t$\xspace}
\newcommand{\bphi}{\boldsymbol{\varphi}}
\newcommand{\bX}{\mathbf{X}}
\newcommand{\bY}{\mathbf{Y}}
\newcommand{\bZ}{\mathcal{Z}}
\newcommand{\ind}{\mathbf{1}}
\newcommand{\Norm}{\mathcal{N}}
\newcommand{\Unif}{\mathrm{Unif}}
\newcommand{\given}{\,\vert\,}
\DeclareRobustCommand{\code}[1]{\texttt{#1}}
\newcommand{\eref}[1]{(\ref{#1})}
\newcommand{\sref}[1]{Section~\ref{#1}}

% ---- Make the math shorthands safe inside PDF bookmarks ----
\pdfstringdefDisableCommands{%
  \def\bphi{phi}%
  \def\boldsymbol#1{#1}%
  \def\bm#1{#1}%
  \def\code#1{#1}%
  \def\skewt{skew-t}%
}

% ---- Theorem environments ----
\theoremstyle{plain}
\newtheorem{proposition}{Proposition}[section]
\newtheorem{lemma}[proposition]{Lemma}
\theoremstyle{definition}
\newtheorem{definition}[proposition]{Definition}
\theoremstyle{remark}
\newtheorem{remark}[proposition]{Remark}

\lstset{
    basicstyle=\ttfamily\footnotesize,
    keywordstyle=\color{blue!60!black},
    commentstyle=\color{gray}\itshape,
    stringstyle=\color{teal},
    showstringspaces=false,
    breaklines=true,
    frame=single,
    framesep=4pt,
    language=R
}

\title{The $\bphi$ Update in \code{mythesis.tex} versus \code{code/R/ar2}:\\
       A Step-by-Step Correspondence Audit}
\author{Technical Note --- \code{thesis\_phi\_update\_audit}}
\date{2026-08-01}

\begin{document}
\maketitle

\begin{abstract}
\noindent
This note checks, step by step, whether the thesis description of how the AR
coefficients $\bphi$ enter and are updated in the MCMC matches
\code{code/R/ar2}. The audit separates three questions that the phrase ``the
$\bphi$ update'' can mean, because the thesis answers them very differently.
\emph{(A)} How $\bphi$ enters the Metropolis--Hastings update of the
time-varying parameters --- Appendix~C of the thesis
(\code{app:full-cond}); the code matches it exactly, including the
head-of-series and end-of-series edge cases. \emph{(B)} How $\bphi$ itself is
drawn --- the code runs a component-wise random-walk Metropolis kernel with
silent rejection on the stationarity triangle, targeting the \emph{exact}
stationary likelihood; the thesis does not describe this block at all.
\emph{(C)} The prior on $\bphi$ --- the thesis states
$\Unif(\text{stationarity triangle})$, whereas the code uses independent
$\Norm(0, 0.5^2)$ factors truncated to that triangle. (C) is a contradiction
between text and code, (B) is an omission, (A) is correct. Two minor prose
inaccuracies in Appendix~C are also recorded. \sref{sec:edits} gives
paste-ready replacement text for \code{mythesis.tex}.
\end{abstract}

\tableofcontents

% ====================================================================
\section{Scope, and a notation collision that must be cleared first}
\label{sec:scope}

\subsection{Two different objects are both called $Y$}

The thesis uses the symbol $Y$ for two things.

\begin{itemize}[nosep]
  \item In the model, $Y_t(\bm{s})$ is the \emph{response} at site $\bm{s}$ on
        day $t$ (thesis \code{eq:fullmodel}), and $\bY_t$ is its
        $n$-vector over sites.
  \item In the thesis's \code{sec:ar2} and throughout Appendix~C
        (\code{app:full-cond}), $\{Y_t\}$ denotes ``any scalar component of
        $\bm{w}^{*}_{tk}$, or $z^{*}_{tk}$ or $\sigma^{2*}_{tk}$'' --- that
        is, one \emph{time-varying parameter on the transformed
        (standard-normal) scale}, at a fixed knot $k$.
\end{itemize}

Both meanings appear inside the same display, thesis \code{eq:full-cond},
whose likelihood factor is written $L(\bY \given Y_t)$: bold $\bY$ is the
data, scalar $Y_t$ is the time-varying parameter. The companion note
\code{tex/transform\_ar\_construction/} avoids the collision by writing
$X^{*}_t$ for the transformed time-varying parameter. \textbf{This note adopts
$X^{*}_t$ throughout}, reserving $\bY_t$ for the data.

\begin{definition}[The three transformed channels]\label{def:channels}
Fix a knot $k \in \{1,\dots,K\}$ and let $X^{*}_{1:n_t}$ denote any one of the
scalar series
\begin{equation}\label{eq:channels}
  X^{*}_{tk} \in
  \bigl\{\, w^{(1)*}_{tk},\; w^{(2)*}_{tk},\; z^{*}_{tk},\; \sigma^{2*}_{tk} \,\bigr\},
\end{equation}
obtained from the probability integral transforms
\code{eq:wstar}--\code{eq:sigmastar} of the thesis. Each channel carries its
own coefficient pair, $\bphi^{(w)}, \bphi^{(z)}, \bphi^{(\sigma)}$, shared
across knots $k$ (and, for $w$, across both coordinates $j = 1,2$).
\end{definition}

\subsection{Three distinct questions}

Once the notation is fixed, ``does the thesis description of the $\bphi$
update match the code?'' splits into three questions, answered in
\crefrange{sec:layerA}{sec:layerC} and summarised in \cref{tab:verdict}.

\begin{table}[htbp]
\centering
\caption{Audit verdict. Thesis labels refer to \code{myLatex/thesis/mythesis.tex};
code paths are relative to \code{code/R/ar2/}.}
\label{tab:verdict}
\small
\begin{tabular}{@{}p{0.055\textwidth} p{0.31\textwidth} p{0.24\textwidth} p{0.30\textwidth}@{}}
\toprule
& \textbf{Question} & \textbf{Thesis} & \textbf{Verdict} \\
\midrule
(A) & How does $\bphi$ enter the MH update of $X^{*}_t$?
    & \code{app:full-cond}, \code{app:ar2-moments}
    & \textbf{Matches}, equation by equation (\sref{sec:layerA}) \\
\addlinespace
(B) & How is $\bphi$ itself updated?
    & --- (nothing)
    & \textbf{Undocumented}; four features missing (\sref{sec:layerB}) \\
\addlinespace
(C) & What is the prior on $\bphi$?
    & \code{eq:stage3-priors}: $\Unif(S_2)$
    & \textbf{Contradicts} the code, which uses truncated
      $\Norm(0,0.5^2)$ (\sref{sec:layerC}) \\
\bottomrule
\end{tabular}
\end{table}

% ====================================================================
\section{The object $\bphi$ parameterises}
\label{sec:prior-object}

\begin{definition}[Standardised stationary AR(2)]\label{def:ar2}
For each channel of \cref{def:channels}, the prior is
\begin{equation}\label{eq:ar2-rec}
  X^{*}_t \;=\; \varphi_1 X^{*}_{t-1} \;+\; \varphi_2 X^{*}_{t-2} \;+\; \epsilon_t,
  \qquad \epsilon_t \overset{\text{iid}}{\sim} \Norm(0, \sigma_\epsilon^2),
  \qquad t \ge 3,
\end{equation}
subject to the unit-variance normalisation $\gamma_0 := \operatorname{Var}(X^{*}_t) = 1$,
which the back-transform requires at \emph{every} $t$, and to stationarity,
\begin{equation}\label{eq:triangle}
  \bphi = (\varphi_1, \varphi_2) \in S_2
  \;:=\;
  \bigl\{ |\varphi_2| < 1,\;\; \varphi_1 + \varphi_2 < 1,\;\; \varphi_2 - \varphi_1 < 1 \bigr\},
\end{equation}
the open triangle with vertices $(-2,-1)$, $(2,-1)$, $(0,1)$.
\end{definition}

\begin{proposition}[Yule--Walker moments]\label{prop:yw}
Under \cref{def:ar2},
\begin{equation}\label{eq:yw}
  \gamma_0 = 1,
  \qquad
  \gamma_1 = \frac{\varphi_1}{1 - \varphi_2},
  \qquad
  \gamma_2 = \varphi_1 \gamma_1 + \varphi_2,
  \qquad
  \sigma_\epsilon^2 = 1 - \varphi_1 \gamma_1 - \varphi_2 \gamma_2 ,
\end{equation}
all functions of $\bphi$ alone.
\end{proposition}

\noindent
\textbf{Code.} \code{get\_ar2\_standard\_moments}, \code{auxfunctions\_ar2.R}
L9--54; the four assignments at L31--34 are \eref{eq:yw} verbatim.
\textbf{Thesis.} \code{eq:gamma12}--\code{eq:sigeps}, derived in
Appendix~B (\code{app:ar2-moments}). \textbf{Match.} Exact. The stationarity
test \eref{eq:triangle} is \code{check\_ar2\_stability}, L2--7 (with a
tolerance $10^{-8}$ on each strict inequality), matching thesis
\code{eq:stab-triangle}.

\begin{proposition}[Joint factorisation]\label{prop:factor}
Under \cref{def:ar2} the joint prior of $X^{*}_{1:n_t}$ factorises as
$p(X^{*}_{1:n_t}) = p_1(X^{*}_1)\, p_2(X^{*}_2 \given X^{*}_1) \prod_{s=3}^{n_t} p_s(X^{*}_s \given X^{*}_{s-1}, X^{*}_{s-2})$
with
\begin{align}
  p_1(x) &= \Norm\!\bigl(x;\; 0,\; 1\bigr), \label{eq:p1}\\
  p_2(x \given x_1) &= \Norm\!\bigl(x;\; \gamma_1 x_1,\; 1 - \gamma_1^2\bigr), \label{eq:p2}\\
  p_s(x \given x_{s-1}, x_{s-2}) &= \Norm\!\bigl(x;\; \varphi_1 x_{s-1} + \varphi_2 x_{s-2},\; \sigma_\epsilon^2\bigr),
    \qquad s \ge 3. \label{eq:ps}
\end{align}
\end{proposition}

\noindent
\textbf{Code.} \code{get\_ar2\_conditional\_params},
\code{auxfunctions\_ar2.R} L56--92, whose three branches are exactly
\eref{eq:p1} (L66--68), \eref{eq:p2} (L70--85), \eref{eq:ps} (L87--91).
\textbf{Thesis.} \code{eq:joint-prior} together with the initial pair
\code{eq:joint-init}; \eref{eq:p2} is the conditional of that bivariate
normal. \textbf{Match.} Exact.

\begin{remark}[$\bphi$ enters through two channels, not one]\label{rem:two-channels}
This is the structural fact behind everything that follows. In \eref{eq:ps},
$\bphi$ enters through the mean coefficients \emph{and} through
$\sigma_\epsilon^2$. In \eref{eq:p2} it enters only through
$\gamma_1 = \varphi_1/(1-\varphi_2)$. The $t = 2$ factor is therefore
$\bphi$-dependent, and
\begin{equation}\label{eq:p2-neq}
  \Norm\!\bigl(\gamma_1 x_1,\; 1-\gamma_1^2\bigr)
  \;=\;
  \Norm\!\bigl(\varphi_1 x_1,\; \sigma_\epsilon^2\bigr)
  \iff \varphi_2 = 0 .
\end{equation}
Using the right-hand form (the recursion extrapolated backwards with a
fictitious $X^{*}_0 = 0$) breaks $\operatorname{Var}(X^{*}_2) = 1$ and hence
breaks the target marginal after the back-transform. The code uses the
left-hand form (``Fix~1'', \code{auxfunctions\_ar2.R} L74--84). Only $p_1$
is $\bphi$-free.
\end{remark}

% ====================================================================
\section{Layer (A): how $\bphi$ enters the update of $X^{*}_t$}
\label{sec:layerA}

\subsection{What the thesis says}

Appendix~C states that a proposal for $X^{*}_t$ changes three factors of
\cref{prop:factor} --- the self factor, the lag-1 factor and the lag-2 factor
--- so that
\begin{equation}\label{eq:fullcond}
  \pi\bigl(X^{*}_t \given X^{*}_{-t}, \bY\bigr)
  \;\propto\;
  \underbrace{L\bigl(\bY \given X^{*}_t\bigr)}_{\text{spatial likelihood}}
  \cdot
  \underbrace{p_t\bigl(X^{*}_t \given X^{*}_{t-1}, X^{*}_{t-2}\bigr)}_{\text{self}}
  \cdot
  \underbrace{p_{t+1}\bigl(X^{*}_{t+1} \given X^{*}_t, X^{*}_{t-1}\bigr)}_{\text{lag-1}}
  \cdot
  \underbrace{p_{t+2}\bigl(X^{*}_{t+2} \given X^{*}_{t+1}, X^{*}_t\bigr)}_{\text{lag-2}},
\end{equation}
with the convention that factors referencing a non-existent time index are
dropped, and that under a symmetric random-walk proposal
$X^{*(c)}_t \sim \Norm(X^{*}_t, s^2)$ the acceptance probability is
$R = \min\{1, \pi(X^{*(c)}_t \given \cdot)/\pi(X^{*}_t \given \cdot)\}$.

\subsection{Making the edge convention explicit}

Write $\mathcal{S}(t) \subseteq \{1,\dots,n_t\}$ for the set of factor indices
in which $X^{*}_t$ appears. \cref{prop:factor} gives
\begin{equation}\label{eq:Sset}
  \mathcal{S}(t) \;=\; \{t,\, t+1,\, t+2\} \cap \{1, \dots, n_t\},
\end{equation}
which unpacks as in \cref{tab:factors}. The acceptance ratio actually
evaluated is then
\begin{equation}\label{eq:logR-latent}
  \log R
  \;=\;
  \underbrace{\log L\bigl(\bY \given X^{*(c)}_t\bigr) - \log L\bigl(\bY \given X^{*}_t\bigr)}_{\text{likelihood part}}
  \;+\;
  \sum_{s \in \mathcal{S}(t)}
  \Bigl[ \log p_s\bigl(\,\cdot \given \cdot\,\bigr)\Big|_{X^{*}_t = X^{*(c)}_t}
       - \log p_s\bigl(\,\cdot \given \cdot\,\bigr)\Big|_{X^{*}_t} \Bigr].
\end{equation}

\begin{table}[htbp]
\centering
\caption{The factor set $\mathcal{S}(t)$ of \eref{eq:Sset}, written out. The
$t = 1$ and $t = 2$ rows are the ``initial-pair'' cases: since
$p(X^{*}_1, X^{*}_2) = p_1(X^{*}_1)\,p_2(X^{*}_2 \given X^{*}_1)$, the head of
the series needs no special code path.}
\label{tab:factors}
\small
\begin{tabular}{@{}l l l@{}}
\toprule
$t$ & Factors containing $X^{*}_t$ & Code guard \\
\midrule
$1$              & $p_1(X^{*}_1)$, \; $p_2(X^{*}_2 \given X^{*}_1)$, \; $p_3(X^{*}_3 \given X^{*}_2, X^{*}_1)$ & self ($t{=}1$ branch), \code{t < nt}, \code{t+2 <= nt} \\
$2$              & $p_2(X^{*}_2 \given X^{*}_1)$, \; $p_3(\cdot)$, \; $p_4(X^{*}_4 \given X^{*}_3, X^{*}_2)$ & self ($t{=}2$ branch), \code{t < nt}, \code{t+2 <= nt} \\
$3 \le t \le n_t - 2$ & $p_t(\cdot)$, \; $p_{t+1}(\cdot)$, \; $p_{t+2}(\cdot)$ & all three \\
$n_t - 1$        & $p_{n_t-1}(\cdot)$, \; $p_{n_t}(\cdot)$ & \code{t+2 <= nt} false \\
$n_t$            & $p_{n_t}(\cdot)$ & \code{t < nt} false \\
\bottomrule
\end{tabular}
\end{table}

\subsection{What the code does}

All three updaters implement \eref{eq:logR-latent} with the guards of
\cref{tab:factors}; see \cref{tab:code-map}. The lag-2 term is the ``Fix~2''
patch documented in \code{tex/where\_we\_are/}, whose
Proposition on detailed balance shows that omitting it makes the chain
target a modified stationary law whose lag-2 autocorrelation misses
$\gamma_2$.

\begin{table}[htbp]
\centering
\caption{Where each factor of \eref{eq:fullcond} lives, by channel. All files
in \code{code/R/ar2/}.}
\label{tab:code-map}
\small
\begin{tabular}{@{}l l l l l@{}}
\toprule
Channel & Updater (file \code{update\_params\_ar2.R}) & self & lag-1 & lag-2 \\
\midrule
$\sigma^{2*}$, $K = 1$ & \code{updateTauTS\_AR2}   & L51--59   & L62--75   & L82--95 \\
$\sigma^{2*}$, $K > 1$ & \code{updateTauTS\_AR2}   & L142--150 & L153--166 & L172--185 \\
$z^{*}$                & \code{updateZTS\_AR2}     & L392--400 & L403--416 & L421--434 \\
$\bm{w}^{*}$           & \code{updateKnotsTS\_AR2} & L548--555 & L566--579 & L584--597 \\
\bottomrule
\end{tabular}
\end{table}

\begin{proposition}[Layer (A) verdict]\label{prop:verdictA}
For every channel and every $t$, the set of prior factors entered into the
acceptance ratio by the code equals $\mathcal{S}(t)$ of \eref{eq:Sset}, and
each factor is evaluated with the branch of \cref{prop:factor} matching its
own time index. The code therefore implements \eref{eq:fullcond} exactly as
Appendix~C describes it.
\end{proposition}

\begin{proof}[Verification]
Take \code{updateZTS\_AR2} as representative. The self term (L392--400) calls
\code{ar2\_conditional\_log\_density} with time index $t$ and lags
\begin{lstlisting}
lag1 <- if (t >= 2) z.star[k, t - 1] else NULL
lag2 <- if (t >= 3) z.star[k, t - 2] else NULL
\end{lstlisting}
so that by
\cref{prop:factor} this selects \eref{eq:p1} at $t=1$, \eref{eq:p2} at $t=2$,
\eref{eq:ps} otherwise. The lag-1 term (L403--416) is guarded by
\code{t < nt} and passes time index $t+1$ with the candidate as its
\emph{lag-1} argument and \code{z.star[k, t-1]} as its lag-2 argument; at
$t = 1$ that argument is \code{NULL} and the $t{+}1 = 2$ branch \eref{eq:p2}
is selected, which is why the head of the series needs no special case. The
lag-2 term (L421--434) is guarded by \code{t+2 <= nt} and passes time index
$t+2$ with \code{z.star[k, t+1]} as lag-1 and the candidate as lag-2. Reading
off the three time indices gives exactly the rows of \cref{tab:factors}. The
same three blocks appear at the line numbers of \cref{tab:code-map} for the
other channels.
\end{proof}

\subsection{Two prose inaccuracies in Appendix~C}
\label{sec:layerA-nits}

\begin{enumerate}[label=(A\arabic*), leftmargin=*]
\item \textbf{The proposal is not scalar for $\bm{w}^{*}$.} Appendix~C writes
      $X^{*(c)}_t \sim \Norm(X^{*}_t, s^2)$, a single-component move. That is
      what \code{updateTauTS\_AR2} and \code{updateZTS\_AR2} do (one $(k,t)$
      cell at a time). \code{updateKnotsTS\_AR2}, however, proposes all
      $K \times 2$ knot coordinates of day $t$ jointly and accepts or rejects
      them as one block (\code{update\_params\_ar2.R} L523--526, single
      \code{log(runif(1)) < R} at L599--606):
      \begin{equation}\label{eq:block-move}
        \bm{w}^{*(c)}_{tk} = \bm{w}^{*}_{tk} + s_{kt}\,\bm{\varepsilon}_{k},
        \qquad \bm{\varepsilon}_k \overset{\text{iid}}{\sim} \Norm_2(\bm{0}, \bm{I}_2),
        \qquad k = 1, \dots, K,
      \end{equation}
      which is still a symmetric random walk, so \eref{eq:logR-latent} is
      unchanged in form, but it is a block move and the text should say so.
\item \textbf{The adaptation band is $[0.25, 0.50]$, not $0.3$--$0.4$.}
      \code{mhupdate} (\code{auxfunctions.R} L99--113) rescales $s$ by $0.8$
      when the running acceptance rate falls below $0.25$ and by $1.2$ when it
      exceeds $0.50$, every $50$ attempts, and this is applied only while
      \code{iter < burn/2} (\code{mcmc\_ar2.R} L556). The stopping of
      adaptation before the sampling phase is worth stating, since it is what
      makes the adaptive scheme valid.
\end{enumerate}

% ====================================================================
\section{Layer (B): the $\bphi$ block itself, which the thesis omits}
\label{sec:layerB}

Appendix~C ends at \eref{eq:fullcond}. Nowhere does the thesis say how
$\bphi^{(w)}, \bphi^{(z)}, \bphi^{(\sigma)}$ are drawn. The code does it in
\code{updatePhiAR2TS} (\code{update\_params\_ar2.R} L225--340), called once
per sweep from each channel's updater (L201--205, L449--453, L613--617). Four
features need documenting.

\subsection{(B1) The target}

Conditionally on the transformed latents, $\bphi$ is independent of the data
and of all other parameters, so its full conditional involves only the AR
prior and its own prior:
\begin{equation}\label{eq:phi-target}
  \pi\bigl(\bphi \given X^{*}\bigr)
  \;\propto\;
  \Biggl[\, \prod_{k=1}^{K} p\bigl(X^{*}_{1:n_t,k} \given \bphi\bigr) \Biggr]\, p(\bphi)\, \ind\{\bphi \in S_2\},
\end{equation}
the product running over knots (and, for $\bm{w}^{*}$, over both coordinates),
which are independent replicates of the same AR(2) law.

\subsection{(B2) The likelihood is the exact stationary one}

Expanding $p(X^{*}_{1:n_t,k} \given \bphi)$ with \cref{prop:factor} and
dropping the $\bphi$-free factor $p_1$, the log-likelihood used in the
acceptance ratio is
\begin{equation}\label{eq:exact-ll}
  \ell(\bphi)
  \;=\;
  \underbrace{\sum_{k} \log \Norm\!\Bigl(X^{*}_{2k};\; \gamma_1 X^{*}_{1k},\; 1 - \gamma_1^2 \Bigr)}_{\text{initial block, } \bphi\text{-dependent}}
  \;+\;
  \underbrace{\sum_{k} \sum_{t=3}^{n_t} \log \Norm\!\Bigl(X^{*}_{tk};\; \varphi_1 X^{*}_{t-1,k} + \varphi_2 X^{*}_{t-2,k},\; \sigma_\epsilon^2 \Bigr)}_{\text{transitions (Box--Jenkins conditional likelihood)}} .
\end{equation}
\textbf{Code.} \code{ar2\_stationary\_loglik}, \code{auxfunctions\_ar2.R}
L122--133; the two summands are L131 and L132 respectively.

\begin{remark}[Why the initial block cannot be dropped, and why AR(1) is exempt]
\label{rem:exact}
Keeping only the transition sum is the Box--Jenkins conditional likelihood.
For AR(1) that is exact up to a $\bphi$-free constant, because the only
discarded factor is $p_1 = \Norm(0,1)$, which does not involve $\varphi$. For
AR(2) the discarded factor is $p_1 \cdot p_2$, and by
\cref{rem:two-channels} $p_2$ \emph{does} involve $\bphi$ through $\gamma_1$.
Dropping it makes the $\bphi$ block and the $X^{*}$ block target different
joint laws, i.e.\ the two Gibbs blocks become incompatible; see
\code{tex/ar2\_phi\_exact\_likelihood/}. Appendix~C already argues the dual
statement for the $X^{*}$ block (``all three prior factors must appear''); the
same $\gamma_1$ dependence produces this second trap, and the thesis states
only the first half.
\end{remark}

\subsection{(B3) The kernel: component-wise random walk with silent rejection}

One sweep updates $\varphi_1$ and then $\varphi_2$, the second step
conditioning on the just-updated first. For $\varphi_1$ (the $\varphi_2$ step
is identical with the roles exchanged):
\begin{align}
  \varphi_1^{(c)} &\sim \Norm\bigl(\varphi_1, s_1^2\bigr) \quad\text{(untruncated)}, \label{eq:phi-prop}\\[2pt]
  \log R &= \ell\bigl(\varphi_1^{(c)}, \varphi_2\bigr) - \ell\bigl(\varphi_1, \varphi_2\bigr)
            + \log p\bigl(\varphi_1^{(c)}\bigr) - \log p\bigl(\varphi_1\bigr),
            \label{eq:phi-ratio}\\[2pt]
  \varphi_1^{\text{new}} &=
  \begin{cases}
    \varphi_1^{(c)}, & \text{if } (\varphi_1^{(c)}, \varphi_2) \in S_2 \text{ and } \log U < \log R, \quad U \sim \Unif(0,1),\\[2pt]
    \varphi_1,       & \text{otherwise (in particular, whenever the candidate leaves } S_2).
  \end{cases}
  \label{eq:silent}
\end{align}
\textbf{Code.} L281--305 ($\varphi_1$) and L308--331 ($\varphi_2$); the
stationarity test at L283/L310 is the $\ind\{\cdot \in S_2\}$ of
\eref{eq:silent}; the moment recomputation is wrapped in \code{tryCatch} so
that a numerically degenerate candidate is also silently rejected.

\begin{lemma}[Silent rejection needs no Hastings correction]\label{lem:silent}
Let $q$ be symmetric, $q(\bphi' \given \bphi) = q(\bphi \given \bphi')$, and
let the kernel accept only candidates in $S_2$, with probability
$\alpha(\bphi, \bphi') = \min\{1, \pi(\bphi')/\pi(\bphi)\}$, otherwise
remaining at $\bphi$. Then for $\bphi, \bphi' \in S_2$,
\begin{equation}\label{eq:db}
  \pi(\bphi)\, q(\bphi' \given \bphi)\, \alpha(\bphi, \bphi')
  \;=\;
  \min\bigl\{ \pi(\bphi), \pi(\bphi') \bigr\}\, q(\bphi' \given \bphi)
  \;=\;
  \pi(\bphi')\, q(\bphi \given \bphi')\, \alpha(\bphi', \bphi),
\end{equation}
so detailed balance holds with respect to $\pi$ restricted to $S_2$; mass
proposed outside $S_2$ contributes only to the holding probability. No
Jacobian or truncation ratio enters.
\end{lemma}

\begin{remark}[Do not ``fix'' this to match the AR(1) code]\label{rem:no-fix}
The AR(1) reference updater \code{updatePhiTS}
(\code{update\_params.R} L1156--1202) instead draws from a normal truncated to
$(-1,1)$ and carries the corresponding asymmetry correction, terms
\code{curlog.cand - canlog.cand} at L1192. The audit in
\code{tex/where\_we\_are/} records that an attempt to impose the analogous
correction on the AR(2) silent-rejection kernel biased the chain toward the
boundary of $S_2$ and was reverted. Both kernels are valid; they are simply
different, and the thesis should describe the one that is used.
\end{remark}

\subsection{(B4) Adaptation and bookkeeping}

The proposal scales $s_1, s_2$ are tuned by the same \code{mhupdate} rule as
in \sref{sec:layerA-nits}, separately for each channel
(\code{mcmc\_ar2.R} L567--581 for $\bm{w}$, and the matching blocks for
$\sigma^2$ and $z$), and only while \code{iter < burn/2}. Attempt counters are
incremented before the stationarity test (L229--230), so an out-of-triangle
proposal counts as a rejection --- which is the correct accounting for the
kernel of \cref{lem:silent}, since such a move genuinely fails.

% ====================================================================
\section{Layer (C): the prior, where text and code disagree}
\label{sec:layerC}

\subsection{The two statements side by side}

The thesis, \code{eq:stage3-priors}, states
\begin{equation}\label{eq:prior-thesis}
  \bigl(\varphi^{(\cdot)}_1,\, \varphi^{(\cdot)}_2\bigr)
  \;\sim\; \Unif(S_2),
  \qquad\text{i.e.}\qquad
  p_{\text{thesis}}(\bphi) \;\propto\; \ind\{\bphi \in S_2\}.
\end{equation}
The code, through the defaults \code{prior.mean = 0}, \code{prior.sd = 0.5} of
\code{updatePhiAR2TS} (L228) used at L295--297 and L322--324, implements
\begin{equation}\label{eq:prior-code}
  p_{\text{code}}(\bphi)
  \;\propto\;
  \Norm\bigl(\varphi_1;\, 0,\, 0.5^2\bigr)\,
  \Norm\bigl(\varphi_2;\, 0,\, 0.5^2\bigr)\,
  \ind\{\bphi \in S_2\},
\end{equation}
the indicator being realised implicitly by \eref{eq:silent}. The arguments
\code{prior.mean} and \code{prior.sd} are never supplied by any production
caller --- \code{updateTauTS\_AR2}, \code{updateZTS\_AR2},
\code{updateKnotsTS\_AR2} all call with the defaults --- and are overridden
only in \code{tests/test\_fix3\_phi\_hastings.R}. So \eref{eq:prior-code} is
effectively hardcoded. The AR(1) baseline is the same story one dimension
down: $\Norm(0, 0.5^2)$ truncated to $(-1,1)$
(\code{update\_params.R} L1190--1191).

\subsection{Why the difference is not cosmetic}

\paragraph{Step 1: the penalty is not flat.}
On $S_2$,
\begin{equation}\label{eq:log-penalty}
  \log p_{\text{code}}(\bphi) - \log p_{\text{code}}(\bm{0})
  \;=\; -\,\frac{\varphi_1^2 + \varphi_2^2}{2 \cdot 0.5^2}
  \;=\; -\,2\bigl(\varphi_1^2 + \varphi_2^2\bigr),
\end{equation}
whereas \eref{eq:prior-thesis} makes this identically $0$. At the generating
values of the simulation study the penalty is $1.53$ nats at
$(\varphi_1,\varphi_2) = (0.80,-0.35)$ and $1.33$ nats at $(0.15, 0.80)$,
relative to the origin. These are small against a likelihood built from
$K(n_t - 1)$ Gaussian factors, which is why
\code{tex/ar2\_phi1\_prior\_support/} measures a posterior shrinkage of at
most $0.04$ and concludes ``no code change''. The point is not that the
sampler is wrong --- it validly targets the posterior induced by
\eref{eq:prior-code} --- but that the thesis reports a prior the study did
not use.

\paragraph{Step 2: the scale is calibrated to the wrong support.}
Under $S_2$ the marginal supports are
\begin{equation}\label{eq:supports}
  \varphi_1 \in (-2, 2), \qquad \varphi_2 \in (-1, 1),
\end{equation}
so $\varphi_1$ has twice the range of the AR(1) coefficient while
\eref{eq:prior-code} keeps the AR(1) scale $0.5$. For the AR(1) pairing the
match is deliberate, $\Phi(1/0.5) - \Phi(-1/0.5) \approx 0.954$, about $95\%$
of the untruncated mass inside the support. For $\varphi_1$ under $S_2$ the
same scale leaves $\Phi(4) - \Phi(-4) > 0.9999$ inside, i.e.\ the prior is far
more informative relative to its support than the AR(1) prior it was ported
from. The mismatch is confined to $\varphi_1$; for $\varphi_2$ support and
scale still match.

\paragraph{Step 3: the two orders are not prior-nested.}
Setting $\varphi_2 = 0$ in \eref{eq:prior-code} recovers the AR(1) prior on
that section. But the AR(2) fit never conditions on $\varphi_2 = 0$; the
comparison in the thesis is against the \emph{marginal} induced on
$\varphi_1$,
\begin{equation}\label{eq:marginal}
  p_{\text{code}}(\varphi_1)
  \;\propto\;
  \Norm\bigl(\varphi_1;\, 0,\, 0.5^2\bigr)
  \int_{\{\varphi_2 :\, (\varphi_1,\varphi_2) \in S_2\}}
  \Norm\bigl(\varphi_2;\, 0,\, 0.5^2\bigr)\, d\varphi_2
  \;\ne\;
  p_{\text{AR(1)}}(\varphi_1),
\end{equation}
because the inner integral depends on $\varphi_1$ through the triangle's
$\varphi_1$-section. Every AR(2)-versus-AR(1) contrast in the thesis
(\code{tab:ar2-sim-paired}, \code{tab:timeblock-brier},
\code{tab:ozone-brier-ci}) therefore carries a small prior confound alongside
the change in order. Recording this caveat in the paper is one of the two
follow-ups left open by \code{tex/ar2\_phi1\_prior\_support/}
(\S Status and decision, 2026-07-15).

% ====================================================================
\section{Recommended edits to \code{mythesis.tex}}
\label{sec:edits}

Nothing in \code{code/R/ar2} needs to change; all three edits are textual.

\paragraph{E1 --- correct the prior (\code{eq:stage3-priors}).}
Replace the $\bphi$ entry by the truncated-normal form and add the
non-nestedness caveat where AR(1) and AR(2) are compared:
\begin{lstlisting}[language=TeX]
(\phi^{(\cdot)}_1,\, \phi^{(\cdot)}_2)
  &\mathrlap{\;\sim\; N(0,\, \sigma^2_\phi \bI_2)\ \text{restricted to the
    stationarity triangle \eref{eq:stab-triangle}},}
% ... and in the hyperparameter sentence: \sigma_\phi = 0.5.
% Caveat sentence, to sit with the AR(1)-vs-AR(2) comparison:
% Because the triangle \eref{eq:stab-triangle} gives $\phi_1$ the support
% $(-2,2)$ while $\phi_2$ retains $(-1,1)$, the common scale
% $\sigma_\phi = 0.5$ is calibrated to the AR(1) support; the induced
% marginal prior on $\phi_1$ under the AR(2) model is therefore not the
% AR(1) prior, and the AR(2)-versus-AR(1) contrasts carry a small prior
% component alongside the change in order. The posterior shrinkage this
% induces was measured at no more than 0.04 across the generating values
% used here.
\end{lstlisting}

\paragraph{E2 --- add the $\bphi$ block to Appendix~C.}
A subsection after the thesis's \code{eq:mh-ratio} covering, in order: the target
\eref{eq:phi-target}; the exact stationary likelihood \eref{eq:exact-ll} with
\cref{rem:exact} on why the initial block survives for AR(2) but is vacuous
for AR(1); the component-wise scan with silent rejection
\eref{eq:phi-prop}--\eref{eq:silent} and \cref{lem:silent}; and one sentence
that a single pair $\bphi^{(\cdot)}$ is shared across knots (and across both
coordinates of $\bm{w}^{*}$).

\paragraph{E3 --- two corrections inside Appendix~C.}
Change ``an acceptance rate near $0.3$--$0.4$'' to the actual rule, adaptation
by factors $0.8/1.2$ whenever the running rate leaves $[0.25, 0.50]$, stopped
after the first half of burn-in; and note that for $\bm{w}^{*}$ the day-$t$
proposal \eref{eq:block-move} is a block move over all $K \times 2$
coordinates rather than the scalar move written for the other two channels.

\paragraph{E4 --- optional, but it is the cause of the confusion.}
Rename the appendix's time-varying parameter from $Y_t$ to $X^{*}_t$ (or
$U_t$), so that \eref{eq:fullcond} does not carry $\bY$ and $Y_t$ with two
different meanings in one display. This matches
\code{tex/transform\_ar\_construction/}.

% ====================================================================
\section{Summary}

\begin{enumerate}[label=\textbf{(\Alph*)}, leftmargin=*]
\item The description of how $\bphi$ enters the update of the time-varying
      parameters --- Appendices~B and~C --- \textbf{matches the code exactly},
      including the Yule--Walker standardisation \eref{eq:yw}, the stationary
      $t = 2$ conditional \eref{eq:p2}, and the three-factor acceptance ratio
      \eref{eq:logR-latent} with its edge conventions. Two prose details need
      correcting (\sref{sec:layerA-nits}).
\item The update of $\bphi$ \emph{itself} is \textbf{not described anywhere in
      the thesis}, although it embodies two deliberate and non-obvious
      decisions: the exact stationary likelihood \eref{eq:exact-ll} and the
      silent-rejection kernel \eref{eq:silent}.
\item The one place the thesis does speak about $\bphi$ directly, its prior,
      \textbf{contradicts the code}: $\Unif(S_2)$ in the text
      \eref{eq:prior-thesis} versus truncated independent $\Norm(0, 0.5^2)$ in
      the implementation \eref{eq:prior-code}, with a documented consequence
      for the AR(1)-versus-AR(2) comparisons \eref{eq:marginal}.
\end{enumerate}

\end{document}
